The recent advancement of Industrial 4.0 has witnessed a significant transformation in modern industrial systems, characterized by increasing complexity, automation, and interconnectedness. To maintain the reliability and safety of these large-scale processes, fault diagnosis has become a critical component for ensuring continuous operations and preventing catastrophic failures \cite{Gao2022, Lei2020}. Data-driven fault diagnosis methods, particularly those based on deep learning, have demonstrated remarkable success in extracting representative features from sensor signals such as vibration, temperature, and pressure \cite{Liu2018}. These methods have fundamentally shifted the paradigm from rule-based and physical modeling to automated feature learning, enabling high-performance diagnostic models in various manufacturing domains. However, the efficacy of these models typically hinges on the availability of massive, high-quality labeled datasets that cover all possible fault conditions. In practice, obtaining such comprehensive data is often infeasible due to the rarity of certain fault modes and the high costs associated with manual labeling in industrial environments.

This data scarcity problem is particularly acute when encountering "unseen" fault types during the testing phase—faults that were never presented during the model's training process. Traditional closed-set fault diagnosis models exhibit a sharp decline in performance when faced with these novel scenarios, as they are restricted to classifying samples into categories predefined by the training labels \cite{Zhang2021}. To overcome this challenge, Zero-Shot Learning (ZSL) has gained traction in the field of intelligent fault diagnosis. ZSL aims to recognize unseen classes by transferring knowledge from seen classes through intermediate semantic representations \cite{Wang2019}. These semantic spaces, often constructed using hand-crafted attributes or word embeddings, serve as a bridge that connects the visual (or sensor-based) features with the label names, allowing the model to generalize to class labels it has never "seen" based on semantic similarity \cite{Lampert2014, Xian2018}.

Non-generative ZSL methods were the first to address fault diagnosis under the zero-shot assumption. These methods typically focus on learning a compatibility function between the feature space and the semantic space or finding a common embedding space where both data modalities can be aligned. For instance, early attempts in fault diagnosis utilized metric learning to measure the distance between sensor features and attribute vectors \cite{Li2020}. Specifically, methods like the Fault Diagnosis with Attribute Transfer (FDAT) and Semantic Consistency Embedding (SCE) paved the way by mapping raw signals into an attribute-based space \cite{Wang2021a}. More advanced architectures, such as Global-Local Attention-based Zero-Shot Learning (GLA-ZSL), leveraged attention mechanisms to focus on discriminative signal segments that correspond to specific fault attributes \cite{Chen2022}. Furthermore, metric-learning-based networks have been optimized to ensure that similar fault types are clustered together in the semantic space, facilitating the classification of unseen categories \cite{Grover2023, Zhang2022}. Despite their efficiency, these non-generative approaches often suffer from the "hubness" problem and the bias towards seen classes, which significantly limits their performance in Generalized Zero-Shot Learning (GZSL) scenarios where the test set contains both seen and unseen samples.

To mitigate the limitations of non-generative models, generative ZSL methods have been introduced, which transform the zero-shot problem into a supervised learning task by synthesizing artificial features for unseen classes. Generative Adversarial Networks (GANs) and Variational Autoencoders (VAEs) are frequently employed to learn the conditional distribution of features given semantic descriptors. Within the context of fault diagnosis, frameworks like Fault-Attribute GAN (FAGAN) and Specialized Robust WGAN (SRWGAN) have been proposed to generate high-fidelity fault features \cite{Yan2021, Zhao2022}. Other researchers have explored hybrid architectures, such as VAEGAN-AR, which combines the distribution-matching capabilities of VAEs with the discriminative power of GANs \cite{Xu2021}. Recent developments have seen the emergence of more sophisticated models like Feature Refinement with Enterprise-driven Embeddings (FREE) and CycleGAN-based Semantic-to-Data (CycleGAN-SD) mapping to enhance the diversity and stability of synthesized data \cite{Devos2023, Kim2022}. Most recently, diffusion models have set new benchmarks in data generation; for example, the Dual-Path Conditional Denoising Diffusion Probabilistic Model with Adaptive Consistency (DP-CDDPM-AC) has demonstrated superior performance in capturing the complex temporal dependencies of industrial signals \cite{Mao2024}. The integration of diffusion models with ZSL has opened new avenues for generating realistic unseen fault features that closely mimic the distribution of real-world data \cite{Ho2020, Rombach2022}.

Despite these advancements, existing ZSL-based diagnostic methods still face three primary limitations. First, the reliance on hand-crafted binary attributes often proves insufficient, as binary 0/1 values cannot capture the subtle, continuous transitions and severity levels of industrial faults, leading to an information bottleneck in the semantic space \cite{Akata2015}. Second, the domain shift problem remains a major hurdle; features synthesized in the latent space based on training distributions often diverge from real-world testing distributions, causing poor generalization to actual unseen samples \cite{Fu2015}. Third, the utilization of the semantic space is often superficial, lacking a deep alignment mechanism that ensures the synthesized features are not only realistic but also semantically consistent with the underlying fault physics \cite{Zhu2019}. These issues collectively hinder the deployment of ZSL models in safety-critical industrial applications where high precision is non-negotiable.

Recent trends in the broader field of artificial intelligence suggest that Large Language Models (LLMs) and multi-modal pre-training models like CLIP can provide a solution to these semantic challenges. In fault diagnosis, there is a burgeoning interest in using natural language descriptions to define fault modes instead of rigid attribute tables. For example, FD-LLM and FaultGPT have explored the use of pre-trained transformers to generate rich, descriptive embeddings from technical manuals and expert reports \cite{Brown2020, Radford2021}. Similarly, applying CLIP-like architectures to structural health monitoring and industrial defect detection has shown that leveraging "language-driven" priors can significantly enhance the robustness of feature representations across different sensors and domains \cite{Li2023a, Wang2023b}. By grounding industrial signals in the rich semantic context of human language, models can achieve a more nuanced understanding of fault relationships that goes beyond mere statistical correlation.

In this paper, we propose a novel framework called Language-Driven Latent Diffusion for Continuous Semantic Alignment (Diff-LM-GZSL) to address the aforementioned challenges. Our method shifts away from discrete attributes, instead utilizing a continuous semantic space derived from pre-trained language models to guide the generation process. We introduce a latent diffusion model that operates on a compressed feature space, allowing for the generation of high-fidelity unseen fault features that are conditioned on complex linguistic descriptions. To further bridge the gap between training and testing distributions, we integrate a semantic alignment mechanism that ensures the generated features remain anchored to their physical meanings. By combining the generative power of diffusion models with the rich context of natural language, Diff-LM-GZSL provides a robust solution for generalized zero-shot fault diagnosis.

The main contributions of this work are three-fold:
\begin{itemize}
    \item We develop a language-driven semantic enhancement module that utilizes LLM/BERT architectures to extract continuous, high-dimensional fault descriptions. This approach replaces conventional hand-crafted binary attributes, providing a more expressive and granular representation of fault characteristics.
    \item We design a Conditional Latent Diffusion Model (CLDM) specialized for the high-fidelity generation of unseen fault features. By operating in the latent space and incorporating continuous semantic guidance, the CLDM reliably synthesizes diverse and precise signal features that represent unseen fault modes.
    \item We propose a semantic alignment mechanism that integrates Maximum Mean Discrepancy (MMD) based distribution matching. This mechanism explicitly mitigates the domain shift between synthesized and real features, ensuring that the diagnostic model remains effective when transferred to real-world industrial datasets.
\end{itemize}

The remainder of this paper is organized as follows: Section II reviews related work in zero-shot learning and fault diagnosis. Section III details the proposed Diff-LM-GZSL framework. Section IV presents the experimental setup and results on three benchmark datasets. Finally, Section V concludes the paper and discusses future research directions.